\documentclass[11pt]{article}
\usepackage[margin=1in]{geometry}
\usepackage{booktabs}
\usepackage{array}

\newcommand{\sym}[1]{\textsuperscript{#1}}

\title{Booktabs Table Stress Case}
\author{Test Corpus}
\date{}

\begin{document}
\maketitle

This document isolates a booktabs table with grouped columns, partial rules, and
significance markers. The table is deliberately compact because the conversion
target should preserve the visual hierarchy created by booktabs rather than
relying on vertical rules.

\begin{table}[htbp]
\centering
\caption{Regression summary with grouped estimates}
\label{tab:booktabs}
\begin{tabular}{lcccc}
\toprule
& \multicolumn{2}{c}{Baseline} & \multicolumn{2}{c}{Adjusted} \\
\cmidrule(lr){2-3}\cmidrule(lr){4-5}
Outcome & Coef. & S.E. & Coef. & S.E. \\
\midrule
Treatment & 0.184\sym{***} & 0.041 & 0.167\sym{***} & 0.039 \\
Post period & 0.072\sym{**} & 0.031 & 0.058\sym{*} & 0.030 \\
Treatment $\times$ Post & 0.119\sym{***} & 0.028 & 0.103\sym{***} & 0.027 \\
Covariate index & -0.033 & 0.024 & -0.029 & 0.023 \\
\midrule
Unit fixed effects & \multicolumn{2}{c}{Yes} & \multicolumn{2}{c}{Yes} \\
Time fixed effects & \multicolumn{2}{c}{No} & \multicolumn{2}{c}{Yes} \\
Observations & \multicolumn{2}{c}{4,812} & \multicolumn{2}{c}{4,812} \\
$R^{2}$ & \multicolumn{2}{c}{0.214} & \multicolumn{2}{c}{0.297} \\
\bottomrule
\end{tabular}
\end{table}

The surrounding prose supplies inline mathematical notation, such as
$\widehat{\beta}=0.103$, to confirm that table conversion can coexist with
ordinary paragraph content.

\end{document}
